% Gerado por .claude/skills/documentar-iniciativa — NÃO EDITE À MÃO.
% A fonte é o markdown do card; regenere após alterá-lo.

\clearpage
\section{Introdução}

\subsection{Contexto}

Este documento consolida os artefatos produzidos pela iniciativa \textbf{Modelo de Precificação NDados} durante o ciclo 2026.1 do núcleo NDados. Reúne o resumo executivo, as atas de benchmark realizadas e os frameworks de discovery derivados delas.

A governança da iniciativa (status, prioridade, checklist e alocação) permanece no Notion. O que está aqui é o registro de campo e a conclusão derivada dele — o material destinado a ser lido daqui a um ano.

\subsection{Por que esta iniciativa existe}

Padronizar esferar de atuação da empresa.

\subsection{Objetivo}

Conceber um modelo matemático de precificação para padronizar o preço de projetos do núcleo.

\subsection{Métricas de sucesso}

faturamento total, taxa de conversão.

\subsection{Como ler este documento}

As atas reproduzem o registro bruto das conversas, separando o que foi dito (seção \emph{Perguntas}) da interpretação da equipe (seção \emph{Síntese}). Os frameworks que vêm depois são análises derivadas: cada um declara, no campo \texttt{origem}, os IDs das atas que o embasam. Nenhuma conclusão apresentada nos frameworks existe sem uma evidência rastreável nas atas.
